\def\StandaloneLibrary{}
% ==============================================================================
% Test Case Reference Library
% Based on ISO/IEC/IEEE 29119-3
% ==============================================================================

% Test Case Attributes
\def\TCID{TC-XXX}
\def\TCTitle{}
\def\TCPurpose{}
\def\TCPriority{}

% Traceability
\def\TCRequirementRef{}
\def\TCDesignRef{}
\def\TCFeatureRef{}

% Preconditions
\def\TCPreconditions{}
\def\TCEnvironment{}
\def\TCAssumptions{}

% Test Data
\def\TCInputData{}
\def\TCInitialState{}

% Execution
\def\TCProcedure{}
\def\TCExpectedResult{}
\def\TCPostConditions{}

% Common Status Values
\def\TCStatusNotRun{Not Executed}
\def\TCStatusPass{Pass}
\def\TCStatusFail{Fail}
\def\TCStatusBlocked{Blocked}
\def\TCStatusInconclusive{Inconclusive}

% Common Priorities
\def\TCPriorityLow{Low}
\def\TCPriorityMedium{Medium}
\def\TCPriorityHigh{High}
\def\TCPriorityCritical{Critical}

% ==============================================================================
% Test Case Template
% ==============================================================================

\newcommand{\TCTemplate}{
	\subsection*{\TCID\ - \TCTitle}
	
	\begin{description}
		
		\item[Purpose]
		\TCPurpose
		
		\item[Priority]
		\TCPriority
		
		\item[Requirement Reference]
		\TCRequirementRef
		
		\item[Design Reference]
		\TCDesignRef
		
		\item[Feature Reference]
		\TCFeatureRef
		
		\item[Preconditions]
		\TCPreconditions
		
		\item[Environment]
		\TCEnvironment
		
		\item[Assumptions]
		\TCAssumptions
		
		\item[Input Data]
		\TCInputData
		
		\item[Initial State]
		\TCInitialState
		
		\item[Procedure] \mbox{}
		\vspace*{-1.5em}
		\TCProcedure
		
		\item[Expected Result]
		\TCExpectedResult
		
		\item[Post Conditions]
		\TCPostConditions
		
	\end{description}
	
	\hrule
	\vspace{1em}
}

% ==============================================================================
% Notes
% ==============================================================================
\ifx\StandaloneLibrary\undefined

\documentclass[letterpaper]{article}

\usepackage[
letterpaper,
margin=1in
]{geometry}

\usepackage[]{embedfile}
\embedfile[
	desc={LaTeX Source}, 
	mimetype=text/plain
]{iso_iec_ieee29119-3_2021.tex}

\pagestyle{empty}

\begin{document}
	{\LARGE\bfseries Jordan's ISO/IEC/IEEE 29119-3 Notes}
	
	\vspace{0.25em}
	
	{\small \today}
	
	\hrule 
	
	% ==============================================================================
	% Section 8.3 Test Case w/ Example
	% ==============================================================================
	\def\TCID{TC-001}

	\def\TCTitle{Power-On Self Test}

	\def\TCPurpose{
		Verify the device successfully completes boot.
	}

	\def\TCPriority{\TCPriorityHigh}

	\def\TCRequirementRef{REQ-001}

	\def\TCPreconditions{
		Device connected to nominal power source.
	}

	\def\TCEnvironment{
		Laboratory ambient conditions.
	}
	
	\def\TCInputData{
		No operator input required.
	}

	\def\TCProcedure{
		\newline
		\begin{enumerate}
			\item Apply power.
			\item Observe startup sequence.
			\item Wait for initialization to complete.
		\end{enumerate}
	}

	\def\TCExpectedResult{
		Device reaches operational state with no errors.
	}

	\def\TCPostConditions{
		Device remains operational.
	}
	\section*{8.3 Test case specification}
	Annex I gives examples of test cases. Here is an example using this library:
	\TCTemplate
\end{document}



\documentclass[letterpaper]{article}
\usepackage[letterpaper,margin=1in]{geometry}

\begin{document}
\def\TCID{TC-0013}
\def\TCTitle{Version-aware package installation and upgrades}
\def\TCPurpose{Verify installed-state discovery, exact architecture and version selectors, configurable prerelease eligibility, SemVer ordering, upgrade validation and auditing, failure recovery, and best-effort multi-package behavior.}
\def\TCPriority{\TCPriorityHigh}
\def\TCRequirementRef{REQ-0013}
\def\TCPreconditions{A built WPM executable is available. The test uses an isolated WPM data-root override and controlled offline repository caches.}
\def\TCEnvironment{Local Windows developer workstation or CI worker.}
\def\TCInputData{Temporary packages covering x86, x64, and any architectures; stable and prerelease SemVer versions; successful and failing install scripts; controlled repository indexes; and an isolated prerelease configuration.}
\def\TCProcedure{\par\begin{enumerate}
\item Build and locally install baseline packages whose metadata creates independent x86, x64, and any installed identities.
\item Configure a controlled repository and seed its offline index and package cache with exact-architecture, any-architecture, stable, prerelease, selected-version, successful, and failing candidates.
\item Upgrade a package having installed x86 and x64 identities without an architecture selector.
\item Exercise named installation and upgrade with \texttt{--arch} and \texttt{--version}, including unavailable, equal, lower, and ZIP-path selector rejection.
\item Read and change the global prerelease setting, add and remove a per-package override, and perform stable and prerelease upgrades.
\item Attempt an any/specific installed-architecture conflict and verify that any candidates do not upgrade specific installed identities.
\item reject an unsigned upgrade before script execution, then explicitly permit it for the positive lifecycle scenarios.
\item Run a multi-package upgrade in which the first package script fails and a later package succeeds.
\item Upgrade WPM itself and verify that a validated candidate executable is launched from cache, waits for the invoking WPM process to exit, and completes installation.
\item Inspect retained archives, staging cleanup, successful and failed upgrade audit records, command summaries, and exit statuses.
\end{enumerate}}
\def\TCExpectedResult{\par\begin{enumerate}
\item Installed identities are discovered from package metadata and the highest retained SemVer is treated as current.
\item x86 and x64 identities are upgraded independently; candidates must exactly match the installed architecture and an any candidate is ignored for a specific identity.
\item Named installation selectors choose the exact requested artifact. Upgrade selectors choose only a strictly newer exact version and never downgrade, reinstall, or silently fall back.
\item An any installation cannot coexist with specific architectures of the same package, and inconsistent state is diagnosed.
\item Prereleases are disabled by default, a package override takes precedence over the global value, stable releases remain eligible, and SemVer prerelease ordering is honored when enabled.
\item Build metadata alone does not cause an upgrade, and invalid repository versions do not prevent selection of valid candidates.
\item An unsigned or otherwise unvalidated candidate is rejected without script execution or successful installation state.
\item Successful upgrades retain both prior and new archives, clean staging, and create an audit record linking old and new versions.
\item A failed script preserves prior state, does not retain the candidate as successfully installed, cleans staging, creates a failure audit, and produces a nonzero result.
\item Multi-package processing continues after a failure, reports each identity as upgraded, current, or failed, and returns nonzero when any identity fails.
\item WPM self-upgrade uses a cached candidate executable to complete replacement after the original process releases the installed executable.
\end{enumerate}}
\def\TCPostConditions{All temporary package sources, archives, repository caches, deployment markers, configuration, audit data, and WPM data directories are removed.}
\TCTemplate
\end{document}
